%% preamble.tex — shared preamble for "Data Engineering" textbook
%% Include with %% preamble.tex — shared preamble for "Data Engineering" textbook
%% Include with %% preamble.tex — shared preamble for "Data Engineering" textbook
%% Include with %% preamble.tex — shared preamble for "Data Engineering" textbook
%% Include with \input{preamble} before \begin{document}
%% Compile: pdflatex → biber → pdflatex → pdflatex

%----------------------------------------------------------------------
% 1. ENCODING AND FONTS
%----------------------------------------------------------------------
\usepackage[T1]{fontenc}
\usepackage[utf8]{inputenc}
\usepackage{newpxtext}          % Palatino-style text (pdfLaTeX compatible)
\usepackage{newpxmath}          % matching math font
\usepackage{inconsolata}        % high-quality monospace for all code
\usepackage{microtype}          % optical margin alignment, letter spacing
\usepackage{setspace}
\setstretch{1.15}               % comfortable reading line spacing

%----------------------------------------------------------------------
% 2. PAGE GEOMETRY
%----------------------------------------------------------------------
\usepackage[
  b5paper,
  top=2.5cm, bottom=2.8cm,
  inner=2.8cm, outer=2.2cm,
  headsep=0.8cm,
  footskip=1.2cm
]{geometry}

%----------------------------------------------------------------------
% 3. COLOR PALETTE
%----------------------------------------------------------------------
\usepackage[dvipsnames,svgnames]{xcolor}
% Primary palette
\definecolor{DEBlue}{RGB}{0,60,120}
\definecolor{DEAccent}{RGB}{210,120,0}
\definecolor{DEGreen}{RGB}{0,110,70}
\definecolor{DERed}{RGB}{170,30,30}
\definecolor{DEGray}{RGB}{80,80,80}
\definecolor{DEPurple}{RGB}{80,0,120}
% Light tints (box backgrounds)
\definecolor{DELightBlue}{RGB}{228,240,255}
\definecolor{DELightAccent}{RGB}{255,245,218}
\definecolor{DELightGreen}{RGB}{218,245,228}
\definecolor{DELightRed}{RGB}{255,232,232}
\definecolor{DELightGray}{RGB}{248,248,248}
\definecolor{DELightPurple}{RGB}{244,235,255}
% Mid tints for lines and rules
\definecolor{DEBlueMid}{RGB}{0,60,120}

%----------------------------------------------------------------------
% 4. MATHEMATICS
%----------------------------------------------------------------------
\usepackage{amsmath}
\usepackage{amssymb}
% newpxmath defines \openbox; amsthm also defines it — clear it first
\let\openbox\relax
\usepackage{amsthm}

%----------------------------------------------------------------------
% 5. TABLES
%----------------------------------------------------------------------
\usepackage{booktabs}
\usepackage{tabularx}
\usepackage{longtable}
\usepackage{array}
\usepackage{multirow}
\usepackage{multicol}
% Convenience column types
\newcolumntype{L}[1]{>{\raggedright\arraybackslash}p{#1}}
\newcolumntype{C}[1]{>{\centering\arraybackslash}p{#1}}
\newcolumntype{R}[1]{>{\raggedleft\arraybackslash}p{#1}}

%----------------------------------------------------------------------
% 6. GRAPHICS AND TikZ
%----------------------------------------------------------------------
\usepackage{graphicx}
\usepackage{tikz}
\usetikzlibrary{
  shapes.geometric,
  shapes.misc,
  shapes.arrows,
  positioning,
  arrows.meta,
  fit,
  backgrounds,
  decorations.pathreplacing,
  decorations.text,
  calc,
  matrix,
  patterns,
  shadows.blur
}
\usepackage{pgfplots}
\pgfplotsset{compat=1.17}

%----------------------------------------------------------------------
% 7. CODE LISTINGS
%----------------------------------------------------------------------
\usepackage{listings}
\renewcommand{\lstlistingname}{Listing}

% PySpark / Python
\lstdefinestyle{pyspark}{
  language=Python,
  basicstyle=\ttfamily\footnotesize,
  backgroundcolor=\color{DEBlue!5},
  frame=single,
  framerule=0.5pt,
  rulecolor=\color{DEBlue!30},
  numbers=left,
  numberstyle=\tiny\color{DEGray},
  numbersep=8pt,
  stepnumber=1,
  breaklines=true,
  breakatwhitespace=false,
  showstringspaces=false,
  keywordstyle=\color{DEBlue}\bfseries,
  commentstyle=\color{DEGreen!80!black}\itshape,
  stringstyle=\color{DERed!80},
  emphstyle=\color{DEAccent}\bfseries,
  emph={SparkSession,SparkContext,RDD,DataFrame,sc,spark},
  morekeywords={as,lambda,True,False,None,self,yield,with},
  captionpos=b,
  xleftmargin=14pt,
  xrightmargin=4pt,
  aboveskip=8pt,
  belowskip=4pt,
  literate={->}{{\texttt{->}}}{2}
}

% HiveQL / SQL
\lstdefinestyle{hiveql}{
  language=SQL,
  basicstyle=\ttfamily\footnotesize,
  backgroundcolor=\color{DEAccent!6},
  frame=single,
  framerule=0.5pt,
  rulecolor=\color{DEAccent!40},
  numbers=left,
  numberstyle=\tiny\color{DEGray},
  numbersep=8pt,
  breaklines=true,
  showstringspaces=false,
  keywordstyle=\color{DEBlue}\bfseries,
  morekeywords={PARTITIONED,STORED,OVERWRITE,ORC,PARQUET,AVRO,
                LOCATION,MSCK,REPAIR,EXTERNAL,TBLPROPERTIES,
                CLUSTERED,SORTED,BUCKETS,ROW,FORMAT,SERDE,
                INPUTFORMAT,OUTPUTFORMAT,LINES,TERMINATED,
                BY,INTO,SKEWED,EXCHANGE},
  commentstyle=\color{DEGreen!80!black}\itshape,
  stringstyle=\color{DERed!80},
  captionpos=b,
  xleftmargin=14pt,
  aboveskip=8pt,
  belowskip=4pt,
}

% Pseudocode
\lstdefinestyle{pseudocode}{
  basicstyle=\ttfamily\footnotesize,
  backgroundcolor=\color{DEGray!8},
  frame=single,
  framerule=0.5pt,
  rulecolor=\color{DEGray!40},
  numbers=left,
  numberstyle=\tiny\color{DEGray},
  numbersep=8pt,
  breaklines=true,
  showstringspaces=false,
  captionpos=b,
  xleftmargin=14pt,
  aboveskip=8pt,
  belowskip=4pt,
}

% Shell / terminal
\lstdefinestyle{shell}{
  basicstyle=\ttfamily\footnotesize\color{white},
  backgroundcolor=\color{black!88},
  frame=single,
  framerule=0.5pt,
  rulecolor=\color{black},
  numbers=none,
  breaklines=true,
  showstringspaces=false,
  captionpos=b,
  xleftmargin=14pt,
  aboveskip=8pt,
  belowskip=4pt,
}

%----------------------------------------------------------------------
% 8. TCOLORBOX ENVIRONMENTS
%----------------------------------------------------------------------
\usepackage[most,breakable]{tcolorbox}
\tcbuselibrary{skins,breakable,listings}

% Research Spotlight (blue)
\newtcolorbox{researchspotlight}[2][]{
  enhanced, breakable,
  colback=DELightBlue,
  colframe=DEBlue,
  colbacktitle=DEBlue,
  coltitle=white,
  fonttitle=\bfseries\small\sffamily,
  title={Research Spotlight\quad\textbar\quad #2},
  left=6pt, right=6pt, top=4pt, bottom=6pt,
  arc=4pt, boxrule=0.8pt,
  attach boxed title to top left={yshift=-2mm, xshift=8pt},
  boxed title style={arc=2pt, colback=DEBlue, size=small},
  #1
}

% Key Insight (amber/accent)
\newtcolorbox{keyinsight}[2][]{
  enhanced, breakable,
  colback=DELightAccent,
  colframe=DEAccent,
  colbacktitle=DEAccent,
  coltitle=white,
  fonttitle=\bfseries\small\sffamily,
  title={Key Insight\quad\textbar\quad #2},
  left=6pt, right=6pt, top=4pt, bottom=6pt,
  arc=4pt, boxrule=0.8pt,
  attach boxed title to top left={yshift=-2mm, xshift=8pt},
  boxed title style={arc=2pt, colback=DEAccent, size=small},
  #1
}

% System Design Note (green)
\newtcolorbox{systemdesign}[2][]{
  enhanced, breakable,
  colback=DELightGreen,
  colframe=DEGreen,
  colbacktitle=DEGreen,
  coltitle=white,
  fonttitle=\bfseries\small\sffamily,
  title={System Design Note\quad\textbar\quad #2},
  left=6pt, right=6pt, top=4pt, bottom=6pt,
  arc=4pt, boxrule=0.8pt,
  attach boxed title to top left={yshift=-2mm, xshift=8pt},
  boxed title style={arc=2pt, colback=DEGreen, size=small},
  #1
}

% Common Misconception (red)
\newtcolorbox{caution}[2][]{
  enhanced, breakable,
  colback=DELightRed,
  colframe=DERed,
  colbacktitle=DERed,
  coltitle=white,
  fonttitle=\bfseries\small\sffamily,
  title={Common Misconception\quad\textbar\quad #2},
  left=6pt, right=6pt, top=4pt, bottom=6pt,
  arc=4pt, boxrule=0.8pt,
  attach boxed title to top left={yshift=-2mm, xshift=8pt},
  boxed title style={arc=2pt, colback=DERed, size=small},
  #1
}

% Production Lesson (gray)
\newtcolorbox{techdebt}[2][]{
  enhanced, breakable,
  colback=DELightGray,
  colframe=DEGray,
  colbacktitle=DEGray,
  coltitle=white,
  fonttitle=\bfseries\small\sffamily,
  title={Production Lesson\quad\textbar\quad #2},
  left=6pt, right=6pt, top=4pt, bottom=6pt,
  arc=4pt, boxrule=0.8pt,
  attach boxed title to top left={yshift=-2mm, xshift=8pt},
  boxed title style={arc=2pt, colback=DEGray, size=small},
  #1
}

% Formal Definition (white + blue border)
\newtcolorbox{defbox}[2][]{
  enhanced, breakable,
  colback=white,
  colframe=DEBlue!55,
  colbacktitle=DEBlue!10,
  coltitle=DEBlue,
  fonttitle=\bfseries\small\sffamily,
  title={Definition\quad\textbar\quad #2},
  left=6pt, right=6pt, top=4pt, bottom=6pt,
  arc=4pt, boxrule=0.8pt,
  #1
}

% Case Study (purple)
\newtcolorbox{casestudy}[2][]{
  enhanced, breakable,
  colback=DELightPurple,
  colframe=DEPurple,
  colbacktitle=DEPurple,
  coltitle=white,
  fonttitle=\bfseries\small\sffamily,
  title={Case Study\quad\textbar\quad #2},
  left=6pt, right=6pt, top=4pt, bottom=6pt,
  arc=4pt, boxrule=0.8pt,
  attach boxed title to top left={yshift=-2mm, xshift=8pt},
  boxed title style={arc=2pt, colback=DEPurple, size=small},
  #1
}

% Chapter Learning Objectives
\newtcolorbox{chapterlearning}[1][]{
  enhanced, breakable,
  colback=DEBlue!5,
  colframe=DEBlue,
  colbacktitle=DEBlue,
  coltitle=white,
  fonttitle=\bfseries\normalsize\sffamily,
  title={Learning Objectives},
  left=8pt, right=8pt, top=6pt, bottom=8pt,
  arc=4pt, boxrule=1pt,
  #1
}

%----------------------------------------------------------------------
% 9. CHAPTER AND SECTION HEADING STYLES
%----------------------------------------------------------------------
\usepackage{titlesec}

% Chapter: oversized gray number, colored title, horizontal rule
\titleformat{\chapter}[display]
  {\normalfont\sffamily}
  {\raggedleft\color{DEBlue!25}\fontsize{90}{90}\selectfont\thechapter}
  {-24pt}
  {\raggedright\color{DEBlue}\fontsize{22}{26}\selectfont\bfseries}
  [\vspace{4pt}\color{DEBlue}\rule{\linewidth}{1.8pt}\vspace{4pt}]

\titlespacing*{\chapter}{0pt}{-40pt}{28pt}

\titleformat{\section}
  {\normalfont\sffamily\Large\bfseries\color{DEBlue}}
  {\thesection}{0.9em}{}
  [\vspace{-4pt}\color{DEBlue!25}\rule{\linewidth}{0.5pt}]

\titleformat{\subsection}
  {\normalfont\sffamily\large\bfseries\color{DEBlue!85}}
  {\thesubsection}{0.9em}{}

\titleformat{\subsubsection}
  {\normalfont\sffamily\normalsize\bfseries\color{DEBlue!75}}
  {\thesubsubsection}{0.9em}{}

%----------------------------------------------------------------------
% 10. HEADERS AND FOOTERS
%----------------------------------------------------------------------
\usepackage{fancyhdr}
\pagestyle{fancy}
\fancyhf{}
\renewcommand{\headrulewidth}{0.4pt}
\renewcommand{\footrulewidth}{0pt}
\fancyhead[LE]{\small\sffamily\color{DEGray}\nouppercase{\leftmark}}
\fancyhead[RO]{\small\sffamily\color{DEGray}\nouppercase{\rightmark}}
\fancyfoot[LE,RO]{\small\sffamily\color{DEGray}\thepage}

\fancypagestyle{plain}{
  \fancyhf{}
  \renewcommand{\headrulewidth}{0pt}
  \fancyfoot[C]{\small\sffamily\color{DEGray}\thepage}
}

%----------------------------------------------------------------------
% 11. CAPTIONS
%----------------------------------------------------------------------
\usepackage{caption}
\usepackage{subcaption}
\captionsetup{
  font=small,
  labelfont={bf,color=DEBlue},
  margin=10pt,
  justification=justified,
}

%----------------------------------------------------------------------
% 12. LISTS
%----------------------------------------------------------------------
\usepackage{enumitem}
\setlist[itemize]{leftmargin=*, topsep=4pt, itemsep=2pt, parsep=0pt}
\setlist[enumerate]{leftmargin=*, topsep=4pt, itemsep=2pt, parsep=0pt}
\setlist[description]{leftmargin=0pt, topsep=4pt, itemsep=2pt}

%----------------------------------------------------------------------
% 13. EPIGRAPH
%----------------------------------------------------------------------
\usepackage{epigraph}
\setlength{\epigraphwidth}{0.62\textwidth}
\renewcommand{\epigraphflush}{flushright}
\renewcommand{\epigraphrule}{0.4pt}

%----------------------------------------------------------------------
% 14. BIBLIOGRAPHY
%----------------------------------------------------------------------
\usepackage[
  style=authoryear,
  backend=biber,
  maxcitenames=3,
  maxbibnames=99,
  giveninits=true,
  uniquename=init,
  dashed=false,
  isbn=false,
  doi=true,
  url=false,
  sorting=nyt
]{biblatex}
\addbibresource{bibliography.bib}

%----------------------------------------------------------------------
% 15. CROSS-REFERENCES AND HYPERLINKS
%----------------------------------------------------------------------
\usepackage[
  colorlinks=true,
  linkcolor=DEBlue,
  citecolor=DEGreen,
  urlcolor=DEAccent,
  pdftitle={Data Engineering: Distributed Systems, Scalable Processing, and Modern Architectures},
  pdfauthor={Rokan Uddin Faruqui},
  pdfsubject={Data Engineering Textbook},
  pdfkeywords={big data, hadoop, spark, data engineering, distributed systems},
  bookmarksnumbered=true,
  bookmarksopen=true,
]{hyperref}
\usepackage[capitalise,noabbrev,nameinlink]{cleveref}

%----------------------------------------------------------------------
% 16. INDEX
%----------------------------------------------------------------------
\usepackage{makeidx}
\makeindex

%----------------------------------------------------------------------
% 17. BLANK PAGES BETWEEN CHAPTERS
%----------------------------------------------------------------------
\usepackage{emptypage}

%----------------------------------------------------------------------
% 18. UTILITY COMMANDS
%----------------------------------------------------------------------
% Defined term (bold + italic) — auto-indexed
\newcommand{\term}[1]{\textbf{\textit{#1}}\index{#1}}
% Software tool name (sans-serif) — auto-indexed
\newcommand{\tool}[1]{\textsf{#1}\index{#1}}
% Inline code
\newcommand{\code}[1]{\texttt{#1}}
% File/dataset name
\newcommand{\dataset}[1]{\texttt{#1}}
% Data-size units (upright, non-breaking)
\newcommand{\KB}{\,\text{KB}}
\newcommand{\MB}{\,\text{MB}}
\newcommand{\GB}{\,\text{GB}}
\newcommand{\TB}{\,\text{TB}}
\newcommand{\PB}{\,\text{PB}}
\newcommand{\EB}{\,\text{EB}}
\newcommand{\ZB}{\,\text{ZB}}

%----------------------------------------------------------------------
% 19. EXERCISE COUNTERS AND MACROS
%----------------------------------------------------------------------
\newcounter{sqctr}[chapter]
\newcounter{apctr}[chapter]
\newcounter{dpctr}[chapter]
\newcounter{pectr}[chapter]

\newcommand{\sq}{\stepcounter{sqctr}%
  \par\vspace{4pt}\noindent
  \textbf{SQ\thechapter.\thesqctr.}\enspace}
\newcommand{\ap}{\stepcounter{apctr}%
  \par\vspace{6pt}\noindent
  \textbf{AP\thechapter.\theapctr.}\enspace}
\newcommand{\dpr}{\stepcounter{dpctr}%
  \par\vspace{6pt}\noindent
  \textbf{DP\thechapter.\thedpctr.}\enspace}
\newcommand{\pe}{\stepcounter{pectr}%
  \par\vspace{6pt}\noindent
  \textbf{PE\thechapter.\thepectr.}\enspace}

% Section heading shorthand for exercise sections
\newcommand{\exercisesection}[1]{%
  \vspace{10pt}%
  {\sffamily\large\bfseries\color{DEBlue} #1}%
  \par\vspace{4pt}%
  {\color{DEBlue!25}\rule{\linewidth}{0.4pt}}%
  \vspace{4pt}%
}
 before \begin{document}
%% Compile: pdflatex → biber → pdflatex → pdflatex

%----------------------------------------------------------------------
% 1. ENCODING AND FONTS
%----------------------------------------------------------------------
\usepackage[T1]{fontenc}
\usepackage[utf8]{inputenc}
\usepackage{newpxtext}          % Palatino-style text (pdfLaTeX compatible)
\usepackage{newpxmath}          % matching math font
\usepackage{inconsolata}        % high-quality monospace for all code
\usepackage{microtype}          % optical margin alignment, letter spacing
\usepackage{setspace}
\setstretch{1.15}               % comfortable reading line spacing

%----------------------------------------------------------------------
% 2. PAGE GEOMETRY
%----------------------------------------------------------------------
\usepackage[
  b5paper,
  top=2.5cm, bottom=2.8cm,
  inner=2.8cm, outer=2.2cm,
  headsep=0.8cm,
  footskip=1.2cm
]{geometry}

%----------------------------------------------------------------------
% 3. COLOR PALETTE
%----------------------------------------------------------------------
\usepackage[dvipsnames,svgnames]{xcolor}
% Primary palette
\definecolor{DEBlue}{RGB}{0,60,120}
\definecolor{DEAccent}{RGB}{210,120,0}
\definecolor{DEGreen}{RGB}{0,110,70}
\definecolor{DERed}{RGB}{170,30,30}
\definecolor{DEGray}{RGB}{80,80,80}
\definecolor{DEPurple}{RGB}{80,0,120}
% Light tints (box backgrounds)
\definecolor{DELightBlue}{RGB}{228,240,255}
\definecolor{DELightAccent}{RGB}{255,245,218}
\definecolor{DELightGreen}{RGB}{218,245,228}
\definecolor{DELightRed}{RGB}{255,232,232}
\definecolor{DELightGray}{RGB}{248,248,248}
\definecolor{DELightPurple}{RGB}{244,235,255}
% Mid tints for lines and rules
\definecolor{DEBlueMid}{RGB}{0,60,120}

%----------------------------------------------------------------------
% 4. MATHEMATICS
%----------------------------------------------------------------------
\usepackage{amsmath}
\usepackage{amssymb}
% newpxmath defines \openbox; amsthm also defines it — clear it first
\let\openbox\relax
\usepackage{amsthm}

%----------------------------------------------------------------------
% 5. TABLES
%----------------------------------------------------------------------
\usepackage{booktabs}
\usepackage{tabularx}
\usepackage{longtable}
\usepackage{array}
\usepackage{multirow}
\usepackage{multicol}
% Convenience column types
\newcolumntype{L}[1]{>{\raggedright\arraybackslash}p{#1}}
\newcolumntype{C}[1]{>{\centering\arraybackslash}p{#1}}
\newcolumntype{R}[1]{>{\raggedleft\arraybackslash}p{#1}}

%----------------------------------------------------------------------
% 6. GRAPHICS AND TikZ
%----------------------------------------------------------------------
\usepackage{graphicx}
\usepackage{tikz}
\usetikzlibrary{
  shapes.geometric,
  shapes.misc,
  shapes.arrows,
  positioning,
  arrows.meta,
  fit,
  backgrounds,
  decorations.pathreplacing,
  decorations.text,
  calc,
  matrix,
  patterns,
  shadows.blur
}
\usepackage{pgfplots}
\pgfplotsset{compat=1.17}

%----------------------------------------------------------------------
% 7. CODE LISTINGS
%----------------------------------------------------------------------
\usepackage{listings}
\renewcommand{\lstlistingname}{Listing}

% PySpark / Python
\lstdefinestyle{pyspark}{
  language=Python,
  basicstyle=\ttfamily\footnotesize,
  backgroundcolor=\color{DEBlue!5},
  frame=single,
  framerule=0.5pt,
  rulecolor=\color{DEBlue!30},
  numbers=left,
  numberstyle=\tiny\color{DEGray},
  numbersep=8pt,
  stepnumber=1,
  breaklines=true,
  breakatwhitespace=false,
  showstringspaces=false,
  keywordstyle=\color{DEBlue}\bfseries,
  commentstyle=\color{DEGreen!80!black}\itshape,
  stringstyle=\color{DERed!80},
  emphstyle=\color{DEAccent}\bfseries,
  emph={SparkSession,SparkContext,RDD,DataFrame,sc,spark},
  morekeywords={as,lambda,True,False,None,self,yield,with},
  captionpos=b,
  xleftmargin=14pt,
  xrightmargin=4pt,
  aboveskip=8pt,
  belowskip=4pt,
  literate={->}{{\texttt{->}}}{2}
}

% HiveQL / SQL
\lstdefinestyle{hiveql}{
  language=SQL,
  basicstyle=\ttfamily\footnotesize,
  backgroundcolor=\color{DEAccent!6},
  frame=single,
  framerule=0.5pt,
  rulecolor=\color{DEAccent!40},
  numbers=left,
  numberstyle=\tiny\color{DEGray},
  numbersep=8pt,
  breaklines=true,
  showstringspaces=false,
  keywordstyle=\color{DEBlue}\bfseries,
  morekeywords={PARTITIONED,STORED,OVERWRITE,ORC,PARQUET,AVRO,
                LOCATION,MSCK,REPAIR,EXTERNAL,TBLPROPERTIES,
                CLUSTERED,SORTED,BUCKETS,ROW,FORMAT,SERDE,
                INPUTFORMAT,OUTPUTFORMAT,LINES,TERMINATED,
                BY,INTO,SKEWED,EXCHANGE},
  commentstyle=\color{DEGreen!80!black}\itshape,
  stringstyle=\color{DERed!80},
  captionpos=b,
  xleftmargin=14pt,
  aboveskip=8pt,
  belowskip=4pt,
}

% Pseudocode
\lstdefinestyle{pseudocode}{
  basicstyle=\ttfamily\footnotesize,
  backgroundcolor=\color{DEGray!8},
  frame=single,
  framerule=0.5pt,
  rulecolor=\color{DEGray!40},
  numbers=left,
  numberstyle=\tiny\color{DEGray},
  numbersep=8pt,
  breaklines=true,
  showstringspaces=false,
  captionpos=b,
  xleftmargin=14pt,
  aboveskip=8pt,
  belowskip=4pt,
}

% Shell / terminal
\lstdefinestyle{shell}{
  basicstyle=\ttfamily\footnotesize\color{white},
  backgroundcolor=\color{black!88},
  frame=single,
  framerule=0.5pt,
  rulecolor=\color{black},
  numbers=none,
  breaklines=true,
  showstringspaces=false,
  captionpos=b,
  xleftmargin=14pt,
  aboveskip=8pt,
  belowskip=4pt,
}

%----------------------------------------------------------------------
% 8. TCOLORBOX ENVIRONMENTS
%----------------------------------------------------------------------
\usepackage[most,breakable]{tcolorbox}
\tcbuselibrary{skins,breakable,listings}

% Research Spotlight (blue)
\newtcolorbox{researchspotlight}[2][]{
  enhanced, breakable,
  colback=DELightBlue,
  colframe=DEBlue,
  colbacktitle=DEBlue,
  coltitle=white,
  fonttitle=\bfseries\small\sffamily,
  title={Research Spotlight\quad\textbar\quad #2},
  left=6pt, right=6pt, top=4pt, bottom=6pt,
  arc=4pt, boxrule=0.8pt,
  attach boxed title to top left={yshift=-2mm, xshift=8pt},
  boxed title style={arc=2pt, colback=DEBlue, size=small},
  #1
}

% Key Insight (amber/accent)
\newtcolorbox{keyinsight}[2][]{
  enhanced, breakable,
  colback=DELightAccent,
  colframe=DEAccent,
  colbacktitle=DEAccent,
  coltitle=white,
  fonttitle=\bfseries\small\sffamily,
  title={Key Insight\quad\textbar\quad #2},
  left=6pt, right=6pt, top=4pt, bottom=6pt,
  arc=4pt, boxrule=0.8pt,
  attach boxed title to top left={yshift=-2mm, xshift=8pt},
  boxed title style={arc=2pt, colback=DEAccent, size=small},
  #1
}

% System Design Note (green)
\newtcolorbox{systemdesign}[2][]{
  enhanced, breakable,
  colback=DELightGreen,
  colframe=DEGreen,
  colbacktitle=DEGreen,
  coltitle=white,
  fonttitle=\bfseries\small\sffamily,
  title={System Design Note\quad\textbar\quad #2},
  left=6pt, right=6pt, top=4pt, bottom=6pt,
  arc=4pt, boxrule=0.8pt,
  attach boxed title to top left={yshift=-2mm, xshift=8pt},
  boxed title style={arc=2pt, colback=DEGreen, size=small},
  #1
}

% Common Misconception (red)
\newtcolorbox{caution}[2][]{
  enhanced, breakable,
  colback=DELightRed,
  colframe=DERed,
  colbacktitle=DERed,
  coltitle=white,
  fonttitle=\bfseries\small\sffamily,
  title={Common Misconception\quad\textbar\quad #2},
  left=6pt, right=6pt, top=4pt, bottom=6pt,
  arc=4pt, boxrule=0.8pt,
  attach boxed title to top left={yshift=-2mm, xshift=8pt},
  boxed title style={arc=2pt, colback=DERed, size=small},
  #1
}

% Production Lesson (gray)
\newtcolorbox{techdebt}[2][]{
  enhanced, breakable,
  colback=DELightGray,
  colframe=DEGray,
  colbacktitle=DEGray,
  coltitle=white,
  fonttitle=\bfseries\small\sffamily,
  title={Production Lesson\quad\textbar\quad #2},
  left=6pt, right=6pt, top=4pt, bottom=6pt,
  arc=4pt, boxrule=0.8pt,
  attach boxed title to top left={yshift=-2mm, xshift=8pt},
  boxed title style={arc=2pt, colback=DEGray, size=small},
  #1
}

% Formal Definition (white + blue border)
\newtcolorbox{defbox}[2][]{
  enhanced, breakable,
  colback=white,
  colframe=DEBlue!55,
  colbacktitle=DEBlue!10,
  coltitle=DEBlue,
  fonttitle=\bfseries\small\sffamily,
  title={Definition\quad\textbar\quad #2},
  left=6pt, right=6pt, top=4pt, bottom=6pt,
  arc=4pt, boxrule=0.8pt,
  #1
}

% Case Study (purple)
\newtcolorbox{casestudy}[2][]{
  enhanced, breakable,
  colback=DELightPurple,
  colframe=DEPurple,
  colbacktitle=DEPurple,
  coltitle=white,
  fonttitle=\bfseries\small\sffamily,
  title={Case Study\quad\textbar\quad #2},
  left=6pt, right=6pt, top=4pt, bottom=6pt,
  arc=4pt, boxrule=0.8pt,
  attach boxed title to top left={yshift=-2mm, xshift=8pt},
  boxed title style={arc=2pt, colback=DEPurple, size=small},
  #1
}

% Chapter Learning Objectives
\newtcolorbox{chapterlearning}[1][]{
  enhanced, breakable,
  colback=DEBlue!5,
  colframe=DEBlue,
  colbacktitle=DEBlue,
  coltitle=white,
  fonttitle=\bfseries\normalsize\sffamily,
  title={Learning Objectives},
  left=8pt, right=8pt, top=6pt, bottom=8pt,
  arc=4pt, boxrule=1pt,
  #1
}

%----------------------------------------------------------------------
% 9. CHAPTER AND SECTION HEADING STYLES
%----------------------------------------------------------------------
\usepackage{titlesec}

% Chapter: oversized gray number, colored title, horizontal rule
\titleformat{\chapter}[display]
  {\normalfont\sffamily}
  {\raggedleft\color{DEBlue!25}\fontsize{90}{90}\selectfont\thechapter}
  {-24pt}
  {\raggedright\color{DEBlue}\fontsize{22}{26}\selectfont\bfseries}
  [\vspace{4pt}\color{DEBlue}\rule{\linewidth}{1.8pt}\vspace{4pt}]

\titlespacing*{\chapter}{0pt}{-40pt}{28pt}

\titleformat{\section}
  {\normalfont\sffamily\Large\bfseries\color{DEBlue}}
  {\thesection}{0.9em}{}
  [\vspace{-4pt}\color{DEBlue!25}\rule{\linewidth}{0.5pt}]

\titleformat{\subsection}
  {\normalfont\sffamily\large\bfseries\color{DEBlue!85}}
  {\thesubsection}{0.9em}{}

\titleformat{\subsubsection}
  {\normalfont\sffamily\normalsize\bfseries\color{DEBlue!75}}
  {\thesubsubsection}{0.9em}{}

%----------------------------------------------------------------------
% 10. HEADERS AND FOOTERS
%----------------------------------------------------------------------
\usepackage{fancyhdr}
\pagestyle{fancy}
\fancyhf{}
\renewcommand{\headrulewidth}{0.4pt}
\renewcommand{\footrulewidth}{0pt}
\fancyhead[LE]{\small\sffamily\color{DEGray}\nouppercase{\leftmark}}
\fancyhead[RO]{\small\sffamily\color{DEGray}\nouppercase{\rightmark}}
\fancyfoot[LE,RO]{\small\sffamily\color{DEGray}\thepage}

\fancypagestyle{plain}{
  \fancyhf{}
  \renewcommand{\headrulewidth}{0pt}
  \fancyfoot[C]{\small\sffamily\color{DEGray}\thepage}
}

%----------------------------------------------------------------------
% 11. CAPTIONS
%----------------------------------------------------------------------
\usepackage{caption}
\usepackage{subcaption}
\captionsetup{
  font=small,
  labelfont={bf,color=DEBlue},
  margin=10pt,
  justification=justified,
}

%----------------------------------------------------------------------
% 12. LISTS
%----------------------------------------------------------------------
\usepackage{enumitem}
\setlist[itemize]{leftmargin=*, topsep=4pt, itemsep=2pt, parsep=0pt}
\setlist[enumerate]{leftmargin=*, topsep=4pt, itemsep=2pt, parsep=0pt}
\setlist[description]{leftmargin=0pt, topsep=4pt, itemsep=2pt}

%----------------------------------------------------------------------
% 13. EPIGRAPH
%----------------------------------------------------------------------
\usepackage{epigraph}
\setlength{\epigraphwidth}{0.62\textwidth}
\renewcommand{\epigraphflush}{flushright}
\renewcommand{\epigraphrule}{0.4pt}

%----------------------------------------------------------------------
% 14. BIBLIOGRAPHY
%----------------------------------------------------------------------
\usepackage[
  style=authoryear,
  backend=biber,
  maxcitenames=3,
  maxbibnames=99,
  giveninits=true,
  uniquename=init,
  dashed=false,
  isbn=false,
  doi=true,
  url=false,
  sorting=nyt
]{biblatex}
\addbibresource{bibliography.bib}

%----------------------------------------------------------------------
% 15. CROSS-REFERENCES AND HYPERLINKS
%----------------------------------------------------------------------
\usepackage[
  colorlinks=true,
  linkcolor=DEBlue,
  citecolor=DEGreen,
  urlcolor=DEAccent,
  pdftitle={Data Engineering: Distributed Systems, Scalable Processing, and Modern Architectures},
  pdfauthor={Rokan Uddin Faruqui},
  pdfsubject={Data Engineering Textbook},
  pdfkeywords={big data, hadoop, spark, data engineering, distributed systems},
  bookmarksnumbered=true,
  bookmarksopen=true,
]{hyperref}
\usepackage[capitalise,noabbrev,nameinlink]{cleveref}

%----------------------------------------------------------------------
% 16. INDEX
%----------------------------------------------------------------------
\usepackage{makeidx}
\makeindex

%----------------------------------------------------------------------
% 17. BLANK PAGES BETWEEN CHAPTERS
%----------------------------------------------------------------------
\usepackage{emptypage}

%----------------------------------------------------------------------
% 18. UTILITY COMMANDS
%----------------------------------------------------------------------
% Defined term (bold + italic) — auto-indexed
\newcommand{\term}[1]{\textbf{\textit{#1}}\index{#1}}
% Software tool name (sans-serif) — auto-indexed
\newcommand{\tool}[1]{\textsf{#1}\index{#1}}
% Inline code
\newcommand{\code}[1]{\texttt{#1}}
% File/dataset name
\newcommand{\dataset}[1]{\texttt{#1}}
% Data-size units (upright, non-breaking)
\newcommand{\KB}{\,\text{KB}}
\newcommand{\MB}{\,\text{MB}}
\newcommand{\GB}{\,\text{GB}}
\newcommand{\TB}{\,\text{TB}}
\newcommand{\PB}{\,\text{PB}}
\newcommand{\EB}{\,\text{EB}}
\newcommand{\ZB}{\,\text{ZB}}

%----------------------------------------------------------------------
% 19. EXERCISE COUNTERS AND MACROS
%----------------------------------------------------------------------
\newcounter{sqctr}[chapter]
\newcounter{apctr}[chapter]
\newcounter{dpctr}[chapter]
\newcounter{pectr}[chapter]

\newcommand{\sq}{\stepcounter{sqctr}%
  \par\vspace{4pt}\noindent
  \textbf{SQ\thechapter.\thesqctr.}\enspace}
\newcommand{\ap}{\stepcounter{apctr}%
  \par\vspace{6pt}\noindent
  \textbf{AP\thechapter.\theapctr.}\enspace}
\newcommand{\dpr}{\stepcounter{dpctr}%
  \par\vspace{6pt}\noindent
  \textbf{DP\thechapter.\thedpctr.}\enspace}
\newcommand{\pe}{\stepcounter{pectr}%
  \par\vspace{6pt}\noindent
  \textbf{PE\thechapter.\thepectr.}\enspace}

% Section heading shorthand for exercise sections
\newcommand{\exercisesection}[1]{%
  \vspace{10pt}%
  {\sffamily\large\bfseries\color{DEBlue} #1}%
  \par\vspace{4pt}%
  {\color{DEBlue!25}\rule{\linewidth}{0.4pt}}%
  \vspace{4pt}%
}
 before \begin{document}
%% Compile: pdflatex → biber → pdflatex → pdflatex

%----------------------------------------------------------------------
% 1. ENCODING AND FONTS
%----------------------------------------------------------------------
\usepackage[T1]{fontenc}
\usepackage[utf8]{inputenc}
\usepackage{newpxtext}          % Palatino-style text (pdfLaTeX compatible)
\usepackage{newpxmath}          % matching math font
\usepackage{inconsolata}        % high-quality monospace for all code
\usepackage{microtype}          % optical margin alignment, letter spacing
\usepackage{setspace}
\setstretch{1.15}               % comfortable reading line spacing

%----------------------------------------------------------------------
% 2. PAGE GEOMETRY
%----------------------------------------------------------------------
\usepackage[
  b5paper,
  top=2.5cm, bottom=2.8cm,
  inner=2.8cm, outer=2.2cm,
  headsep=0.8cm,
  footskip=1.2cm
]{geometry}

%----------------------------------------------------------------------
% 3. COLOR PALETTE
%----------------------------------------------------------------------
\usepackage[dvipsnames,svgnames]{xcolor}
% Primary palette
\definecolor{DEBlue}{RGB}{0,60,120}
\definecolor{DEAccent}{RGB}{210,120,0}
\definecolor{DEGreen}{RGB}{0,110,70}
\definecolor{DERed}{RGB}{170,30,30}
\definecolor{DEGray}{RGB}{80,80,80}
\definecolor{DEPurple}{RGB}{80,0,120}
% Light tints (box backgrounds)
\definecolor{DELightBlue}{RGB}{228,240,255}
\definecolor{DELightAccent}{RGB}{255,245,218}
\definecolor{DELightGreen}{RGB}{218,245,228}
\definecolor{DELightRed}{RGB}{255,232,232}
\definecolor{DELightGray}{RGB}{248,248,248}
\definecolor{DELightPurple}{RGB}{244,235,255}
% Mid tints for lines and rules
\definecolor{DEBlueMid}{RGB}{0,60,120}

%----------------------------------------------------------------------
% 4. MATHEMATICS
%----------------------------------------------------------------------
\usepackage{amsmath}
\usepackage{amssymb}
% newpxmath defines \openbox; amsthm also defines it — clear it first
\let\openbox\relax
\usepackage{amsthm}

%----------------------------------------------------------------------
% 5. TABLES
%----------------------------------------------------------------------
\usepackage{booktabs}
\usepackage{tabularx}
\usepackage{longtable}
\usepackage{array}
\usepackage{multirow}
\usepackage{multicol}
% Convenience column types
\newcolumntype{L}[1]{>{\raggedright\arraybackslash}p{#1}}
\newcolumntype{C}[1]{>{\centering\arraybackslash}p{#1}}
\newcolumntype{R}[1]{>{\raggedleft\arraybackslash}p{#1}}

%----------------------------------------------------------------------
% 6. GRAPHICS AND TikZ
%----------------------------------------------------------------------
\usepackage{graphicx}
\usepackage{tikz}
\usetikzlibrary{
  shapes.geometric,
  shapes.misc,
  shapes.arrows,
  positioning,
  arrows.meta,
  fit,
  backgrounds,
  decorations.pathreplacing,
  decorations.text,
  calc,
  matrix,
  patterns,
  shadows.blur
}
\usepackage{pgfplots}
\pgfplotsset{compat=1.17}

%----------------------------------------------------------------------
% 7. CODE LISTINGS
%----------------------------------------------------------------------
\usepackage{listings}
\renewcommand{\lstlistingname}{Listing}

% PySpark / Python
\lstdefinestyle{pyspark}{
  language=Python,
  basicstyle=\ttfamily\footnotesize,
  backgroundcolor=\color{DEBlue!5},
  frame=single,
  framerule=0.5pt,
  rulecolor=\color{DEBlue!30},
  numbers=left,
  numberstyle=\tiny\color{DEGray},
  numbersep=8pt,
  stepnumber=1,
  breaklines=true,
  breakatwhitespace=false,
  showstringspaces=false,
  keywordstyle=\color{DEBlue}\bfseries,
  commentstyle=\color{DEGreen!80!black}\itshape,
  stringstyle=\color{DERed!80},
  emphstyle=\color{DEAccent}\bfseries,
  emph={SparkSession,SparkContext,RDD,DataFrame,sc,spark},
  morekeywords={as,lambda,True,False,None,self,yield,with},
  captionpos=b,
  xleftmargin=14pt,
  xrightmargin=4pt,
  aboveskip=8pt,
  belowskip=4pt,
  literate={->}{{\texttt{->}}}{2}
}

% HiveQL / SQL
\lstdefinestyle{hiveql}{
  language=SQL,
  basicstyle=\ttfamily\footnotesize,
  backgroundcolor=\color{DEAccent!6},
  frame=single,
  framerule=0.5pt,
  rulecolor=\color{DEAccent!40},
  numbers=left,
  numberstyle=\tiny\color{DEGray},
  numbersep=8pt,
  breaklines=true,
  showstringspaces=false,
  keywordstyle=\color{DEBlue}\bfseries,
  morekeywords={PARTITIONED,STORED,OVERWRITE,ORC,PARQUET,AVRO,
                LOCATION,MSCK,REPAIR,EXTERNAL,TBLPROPERTIES,
                CLUSTERED,SORTED,BUCKETS,ROW,FORMAT,SERDE,
                INPUTFORMAT,OUTPUTFORMAT,LINES,TERMINATED,
                BY,INTO,SKEWED,EXCHANGE},
  commentstyle=\color{DEGreen!80!black}\itshape,
  stringstyle=\color{DERed!80},
  captionpos=b,
  xleftmargin=14pt,
  aboveskip=8pt,
  belowskip=4pt,
}

% Pseudocode
\lstdefinestyle{pseudocode}{
  basicstyle=\ttfamily\footnotesize,
  backgroundcolor=\color{DEGray!8},
  frame=single,
  framerule=0.5pt,
  rulecolor=\color{DEGray!40},
  numbers=left,
  numberstyle=\tiny\color{DEGray},
  numbersep=8pt,
  breaklines=true,
  showstringspaces=false,
  captionpos=b,
  xleftmargin=14pt,
  aboveskip=8pt,
  belowskip=4pt,
}

% Shell / terminal
\lstdefinestyle{shell}{
  basicstyle=\ttfamily\footnotesize\color{white},
  backgroundcolor=\color{black!88},
  frame=single,
  framerule=0.5pt,
  rulecolor=\color{black},
  numbers=none,
  breaklines=true,
  showstringspaces=false,
  captionpos=b,
  xleftmargin=14pt,
  aboveskip=8pt,
  belowskip=4pt,
}

%----------------------------------------------------------------------
% 8. TCOLORBOX ENVIRONMENTS
%----------------------------------------------------------------------
\usepackage[most,breakable]{tcolorbox}
\tcbuselibrary{skins,breakable,listings}

% Research Spotlight (blue)
\newtcolorbox{researchspotlight}[2][]{
  enhanced, breakable,
  colback=DELightBlue,
  colframe=DEBlue,
  colbacktitle=DEBlue,
  coltitle=white,
  fonttitle=\bfseries\small\sffamily,
  title={Research Spotlight\quad\textbar\quad #2},
  left=6pt, right=6pt, top=4pt, bottom=6pt,
  arc=4pt, boxrule=0.8pt,
  attach boxed title to top left={yshift=-2mm, xshift=8pt},
  boxed title style={arc=2pt, colback=DEBlue, size=small},
  #1
}

% Key Insight (amber/accent)
\newtcolorbox{keyinsight}[2][]{
  enhanced, breakable,
  colback=DELightAccent,
  colframe=DEAccent,
  colbacktitle=DEAccent,
  coltitle=white,
  fonttitle=\bfseries\small\sffamily,
  title={Key Insight\quad\textbar\quad #2},
  left=6pt, right=6pt, top=4pt, bottom=6pt,
  arc=4pt, boxrule=0.8pt,
  attach boxed title to top left={yshift=-2mm, xshift=8pt},
  boxed title style={arc=2pt, colback=DEAccent, size=small},
  #1
}

% System Design Note (green)
\newtcolorbox{systemdesign}[2][]{
  enhanced, breakable,
  colback=DELightGreen,
  colframe=DEGreen,
  colbacktitle=DEGreen,
  coltitle=white,
  fonttitle=\bfseries\small\sffamily,
  title={System Design Note\quad\textbar\quad #2},
  left=6pt, right=6pt, top=4pt, bottom=6pt,
  arc=4pt, boxrule=0.8pt,
  attach boxed title to top left={yshift=-2mm, xshift=8pt},
  boxed title style={arc=2pt, colback=DEGreen, size=small},
  #1
}

% Common Misconception (red)
\newtcolorbox{caution}[2][]{
  enhanced, breakable,
  colback=DELightRed,
  colframe=DERed,
  colbacktitle=DERed,
  coltitle=white,
  fonttitle=\bfseries\small\sffamily,
  title={Common Misconception\quad\textbar\quad #2},
  left=6pt, right=6pt, top=4pt, bottom=6pt,
  arc=4pt, boxrule=0.8pt,
  attach boxed title to top left={yshift=-2mm, xshift=8pt},
  boxed title style={arc=2pt, colback=DERed, size=small},
  #1
}

% Production Lesson (gray)
\newtcolorbox{techdebt}[2][]{
  enhanced, breakable,
  colback=DELightGray,
  colframe=DEGray,
  colbacktitle=DEGray,
  coltitle=white,
  fonttitle=\bfseries\small\sffamily,
  title={Production Lesson\quad\textbar\quad #2},
  left=6pt, right=6pt, top=4pt, bottom=6pt,
  arc=4pt, boxrule=0.8pt,
  attach boxed title to top left={yshift=-2mm, xshift=8pt},
  boxed title style={arc=2pt, colback=DEGray, size=small},
  #1
}

% Formal Definition (white + blue border)
\newtcolorbox{defbox}[2][]{
  enhanced, breakable,
  colback=white,
  colframe=DEBlue!55,
  colbacktitle=DEBlue!10,
  coltitle=DEBlue,
  fonttitle=\bfseries\small\sffamily,
  title={Definition\quad\textbar\quad #2},
  left=6pt, right=6pt, top=4pt, bottom=6pt,
  arc=4pt, boxrule=0.8pt,
  #1
}

% Case Study (purple)
\newtcolorbox{casestudy}[2][]{
  enhanced, breakable,
  colback=DELightPurple,
  colframe=DEPurple,
  colbacktitle=DEPurple,
  coltitle=white,
  fonttitle=\bfseries\small\sffamily,
  title={Case Study\quad\textbar\quad #2},
  left=6pt, right=6pt, top=4pt, bottom=6pt,
  arc=4pt, boxrule=0.8pt,
  attach boxed title to top left={yshift=-2mm, xshift=8pt},
  boxed title style={arc=2pt, colback=DEPurple, size=small},
  #1
}

% Chapter Learning Objectives
\newtcolorbox{chapterlearning}[1][]{
  enhanced, breakable,
  colback=DEBlue!5,
  colframe=DEBlue,
  colbacktitle=DEBlue,
  coltitle=white,
  fonttitle=\bfseries\normalsize\sffamily,
  title={Learning Objectives},
  left=8pt, right=8pt, top=6pt, bottom=8pt,
  arc=4pt, boxrule=1pt,
  #1
}

%----------------------------------------------------------------------
% 9. CHAPTER AND SECTION HEADING STYLES
%----------------------------------------------------------------------
\usepackage{titlesec}

% Chapter: oversized gray number, colored title, horizontal rule
\titleformat{\chapter}[display]
  {\normalfont\sffamily}
  {\raggedleft\color{DEBlue!25}\fontsize{90}{90}\selectfont\thechapter}
  {-24pt}
  {\raggedright\color{DEBlue}\fontsize{22}{26}\selectfont\bfseries}
  [\vspace{4pt}\color{DEBlue}\rule{\linewidth}{1.8pt}\vspace{4pt}]

\titlespacing*{\chapter}{0pt}{-40pt}{28pt}

\titleformat{\section}
  {\normalfont\sffamily\Large\bfseries\color{DEBlue}}
  {\thesection}{0.9em}{}
  [\vspace{-4pt}\color{DEBlue!25}\rule{\linewidth}{0.5pt}]

\titleformat{\subsection}
  {\normalfont\sffamily\large\bfseries\color{DEBlue!85}}
  {\thesubsection}{0.9em}{}

\titleformat{\subsubsection}
  {\normalfont\sffamily\normalsize\bfseries\color{DEBlue!75}}
  {\thesubsubsection}{0.9em}{}

%----------------------------------------------------------------------
% 10. HEADERS AND FOOTERS
%----------------------------------------------------------------------
\usepackage{fancyhdr}
\pagestyle{fancy}
\fancyhf{}
\renewcommand{\headrulewidth}{0.4pt}
\renewcommand{\footrulewidth}{0pt}
\fancyhead[LE]{\small\sffamily\color{DEGray}\nouppercase{\leftmark}}
\fancyhead[RO]{\small\sffamily\color{DEGray}\nouppercase{\rightmark}}
\fancyfoot[LE,RO]{\small\sffamily\color{DEGray}\thepage}

\fancypagestyle{plain}{
  \fancyhf{}
  \renewcommand{\headrulewidth}{0pt}
  \fancyfoot[C]{\small\sffamily\color{DEGray}\thepage}
}

%----------------------------------------------------------------------
% 11. CAPTIONS
%----------------------------------------------------------------------
\usepackage{caption}
\usepackage{subcaption}
\captionsetup{
  font=small,
  labelfont={bf,color=DEBlue},
  margin=10pt,
  justification=justified,
}

%----------------------------------------------------------------------
% 12. LISTS
%----------------------------------------------------------------------
\usepackage{enumitem}
\setlist[itemize]{leftmargin=*, topsep=4pt, itemsep=2pt, parsep=0pt}
\setlist[enumerate]{leftmargin=*, topsep=4pt, itemsep=2pt, parsep=0pt}
\setlist[description]{leftmargin=0pt, topsep=4pt, itemsep=2pt}

%----------------------------------------------------------------------
% 13. EPIGRAPH
%----------------------------------------------------------------------
\usepackage{epigraph}
\setlength{\epigraphwidth}{0.62\textwidth}
\renewcommand{\epigraphflush}{flushright}
\renewcommand{\epigraphrule}{0.4pt}

%----------------------------------------------------------------------
% 14. BIBLIOGRAPHY
%----------------------------------------------------------------------
\usepackage[
  style=authoryear,
  backend=biber,
  maxcitenames=3,
  maxbibnames=99,
  giveninits=true,
  uniquename=init,
  dashed=false,
  isbn=false,
  doi=true,
  url=false,
  sorting=nyt
]{biblatex}
\addbibresource{bibliography.bib}

%----------------------------------------------------------------------
% 15. CROSS-REFERENCES AND HYPERLINKS
%----------------------------------------------------------------------
\usepackage[
  colorlinks=true,
  linkcolor=DEBlue,
  citecolor=DEGreen,
  urlcolor=DEAccent,
  pdftitle={Data Engineering: Distributed Systems, Scalable Processing, and Modern Architectures},
  pdfauthor={Rokan Uddin Faruqui},
  pdfsubject={Data Engineering Textbook},
  pdfkeywords={big data, hadoop, spark, data engineering, distributed systems},
  bookmarksnumbered=true,
  bookmarksopen=true,
]{hyperref}
\usepackage[capitalise,noabbrev,nameinlink]{cleveref}

%----------------------------------------------------------------------
% 16. INDEX
%----------------------------------------------------------------------
\usepackage{makeidx}
\makeindex

%----------------------------------------------------------------------
% 17. BLANK PAGES BETWEEN CHAPTERS
%----------------------------------------------------------------------
\usepackage{emptypage}

%----------------------------------------------------------------------
% 18. UTILITY COMMANDS
%----------------------------------------------------------------------
% Defined term (bold + italic) — auto-indexed
\newcommand{\term}[1]{\textbf{\textit{#1}}\index{#1}}
% Software tool name (sans-serif) — auto-indexed
\newcommand{\tool}[1]{\textsf{#1}\index{#1}}
% Inline code
\newcommand{\code}[1]{\texttt{#1}}
% File/dataset name
\newcommand{\dataset}[1]{\texttt{#1}}
% Data-size units (upright, non-breaking)
\newcommand{\KB}{\,\text{KB}}
\newcommand{\MB}{\,\text{MB}}
\newcommand{\GB}{\,\text{GB}}
\newcommand{\TB}{\,\text{TB}}
\newcommand{\PB}{\,\text{PB}}
\newcommand{\EB}{\,\text{EB}}
\newcommand{\ZB}{\,\text{ZB}}

%----------------------------------------------------------------------
% 19. EXERCISE COUNTERS AND MACROS
%----------------------------------------------------------------------
\newcounter{sqctr}[chapter]
\newcounter{apctr}[chapter]
\newcounter{dpctr}[chapter]
\newcounter{pectr}[chapter]

\newcommand{\sq}{\stepcounter{sqctr}%
  \par\vspace{4pt}\noindent
  \textbf{SQ\thechapter.\thesqctr.}\enspace}
\newcommand{\ap}{\stepcounter{apctr}%
  \par\vspace{6pt}\noindent
  \textbf{AP\thechapter.\theapctr.}\enspace}
\newcommand{\dpr}{\stepcounter{dpctr}%
  \par\vspace{6pt}\noindent
  \textbf{DP\thechapter.\thedpctr.}\enspace}
\newcommand{\pe}{\stepcounter{pectr}%
  \par\vspace{6pt}\noindent
  \textbf{PE\thechapter.\thepectr.}\enspace}

% Section heading shorthand for exercise sections
\newcommand{\exercisesection}[1]{%
  \vspace{10pt}%
  {\sffamily\large\bfseries\color{DEBlue} #1}%
  \par\vspace{4pt}%
  {\color{DEBlue!25}\rule{\linewidth}{0.4pt}}%
  \vspace{4pt}%
}
 before \begin{document}
%% Compile: pdflatex → biber → pdflatex → pdflatex

%----------------------------------------------------------------------
% 1. ENCODING AND FONTS
%----------------------------------------------------------------------
\usepackage[T1]{fontenc}
\usepackage[utf8]{inputenc}
\usepackage{newpxtext}          % Palatino-style text (pdfLaTeX compatible)
\usepackage{newpxmath}          % matching math font
\usepackage{inconsolata}        % high-quality monospace for all code
\usepackage{microtype}          % optical margin alignment, letter spacing
\usepackage{setspace}
\setstretch{1.15}               % comfortable reading line spacing

%----------------------------------------------------------------------
% 2. PAGE GEOMETRY
%----------------------------------------------------------------------
\usepackage[
  b5paper,
  top=2.5cm, bottom=2.8cm,
  inner=2.8cm, outer=2.2cm,
  headsep=0.8cm,
  footskip=1.2cm
]{geometry}

%----------------------------------------------------------------------
% 3. COLOR PALETTE
%----------------------------------------------------------------------
\usepackage[dvipsnames,svgnames]{xcolor}
% Primary palette
\definecolor{DEBlue}{RGB}{0,60,120}
\definecolor{DEAccent}{RGB}{210,120,0}
\definecolor{DEGreen}{RGB}{0,110,70}
\definecolor{DERed}{RGB}{170,30,30}
\definecolor{DEGray}{RGB}{80,80,80}
\definecolor{DEPurple}{RGB}{80,0,120}
% Light tints (box backgrounds)
\definecolor{DELightBlue}{RGB}{228,240,255}
\definecolor{DELightAccent}{RGB}{255,245,218}
\definecolor{DELightGreen}{RGB}{218,245,228}
\definecolor{DELightRed}{RGB}{255,232,232}
\definecolor{DELightGray}{RGB}{248,248,248}
\definecolor{DELightPurple}{RGB}{244,235,255}
% Mid tints for lines and rules
\definecolor{DEBlueMid}{RGB}{0,60,120}

%----------------------------------------------------------------------
% 4. MATHEMATICS
%----------------------------------------------------------------------
\usepackage{amsmath}
\usepackage{amssymb}
% newpxmath defines \openbox; amsthm also defines it — clear it first
\let\openbox\relax
\usepackage{amsthm}

%----------------------------------------------------------------------
% 5. TABLES
%----------------------------------------------------------------------
\usepackage{booktabs}
\usepackage{tabularx}
\usepackage{longtable}
\usepackage{array}
\usepackage{multirow}
\usepackage{multicol}
% Convenience column types
\newcolumntype{L}[1]{>{\raggedright\arraybackslash}p{#1}}
\newcolumntype{C}[1]{>{\centering\arraybackslash}p{#1}}
\newcolumntype{R}[1]{>{\raggedleft\arraybackslash}p{#1}}

%----------------------------------------------------------------------
% 6. GRAPHICS AND TikZ
%----------------------------------------------------------------------
\usepackage{graphicx}
\usepackage{tikz}
\usetikzlibrary{
  shapes.geometric,
  shapes.misc,
  shapes.arrows,
  positioning,
  arrows.meta,
  fit,
  backgrounds,
  decorations.pathreplacing,
  decorations.text,
  calc,
  matrix,
  patterns,
  shadows.blur
}
\usepackage{pgfplots}
\pgfplotsset{compat=1.17}

%----------------------------------------------------------------------
% 7. CODE LISTINGS
%----------------------------------------------------------------------
\usepackage{listings}
\renewcommand{\lstlistingname}{Listing}

% PySpark / Python
\lstdefinestyle{pyspark}{
  language=Python,
  basicstyle=\ttfamily\footnotesize,
  backgroundcolor=\color{DEBlue!5},
  frame=single,
  framerule=0.5pt,
  rulecolor=\color{DEBlue!30},
  numbers=left,
  numberstyle=\tiny\color{DEGray},
  numbersep=8pt,
  stepnumber=1,
  breaklines=true,
  breakatwhitespace=false,
  showstringspaces=false,
  keywordstyle=\color{DEBlue}\bfseries,
  commentstyle=\color{DEGreen!80!black}\itshape,
  stringstyle=\color{DERed!80},
  emphstyle=\color{DEAccent}\bfseries,
  emph={SparkSession,SparkContext,RDD,DataFrame,sc,spark},
  morekeywords={as,lambda,True,False,None,self,yield,with},
  captionpos=b,
  xleftmargin=14pt,
  xrightmargin=4pt,
  aboveskip=8pt,
  belowskip=4pt,
  literate={->}{{\texttt{->}}}{2}
}

% HiveQL / SQL
\lstdefinestyle{hiveql}{
  language=SQL,
  basicstyle=\ttfamily\footnotesize,
  backgroundcolor=\color{DEAccent!6},
  frame=single,
  framerule=0.5pt,
  rulecolor=\color{DEAccent!40},
  numbers=left,
  numberstyle=\tiny\color{DEGray},
  numbersep=8pt,
  breaklines=true,
  showstringspaces=false,
  keywordstyle=\color{DEBlue}\bfseries,
  morekeywords={PARTITIONED,STORED,OVERWRITE,ORC,PARQUET,AVRO,
                LOCATION,MSCK,REPAIR,EXTERNAL,TBLPROPERTIES,
                CLUSTERED,SORTED,BUCKETS,ROW,FORMAT,SERDE,
                INPUTFORMAT,OUTPUTFORMAT,LINES,TERMINATED,
                BY,INTO,SKEWED,EXCHANGE},
  commentstyle=\color{DEGreen!80!black}\itshape,
  stringstyle=\color{DERed!80},
  captionpos=b,
  xleftmargin=14pt,
  aboveskip=8pt,
  belowskip=4pt,
}

% Pseudocode
\lstdefinestyle{pseudocode}{
  basicstyle=\ttfamily\footnotesize,
  backgroundcolor=\color{DEGray!8},
  frame=single,
  framerule=0.5pt,
  rulecolor=\color{DEGray!40},
  numbers=left,
  numberstyle=\tiny\color{DEGray},
  numbersep=8pt,
  breaklines=true,
  showstringspaces=false,
  captionpos=b,
  xleftmargin=14pt,
  aboveskip=8pt,
  belowskip=4pt,
}

% Shell / terminal
\lstdefinestyle{shell}{
  basicstyle=\ttfamily\footnotesize\color{white},
  backgroundcolor=\color{black!88},
  frame=single,
  framerule=0.5pt,
  rulecolor=\color{black},
  numbers=none,
  breaklines=true,
  showstringspaces=false,
  captionpos=b,
  xleftmargin=14pt,
  aboveskip=8pt,
  belowskip=4pt,
}

%----------------------------------------------------------------------
% 8. TCOLORBOX ENVIRONMENTS
%----------------------------------------------------------------------
\usepackage[most,breakable]{tcolorbox}
\tcbuselibrary{skins,breakable,listings}

% Research Spotlight (blue)
\newtcolorbox{researchspotlight}[2][]{
  enhanced, breakable,
  colback=DELightBlue,
  colframe=DEBlue,
  colbacktitle=DEBlue,
  coltitle=white,
  fonttitle=\bfseries\small\sffamily,
  title={Research Spotlight\quad\textbar\quad #2},
  left=6pt, right=6pt, top=4pt, bottom=6pt,
  arc=4pt, boxrule=0.8pt,
  attach boxed title to top left={yshift=-2mm, xshift=8pt},
  boxed title style={arc=2pt, colback=DEBlue, size=small},
  #1
}

% Key Insight (amber/accent)
\newtcolorbox{keyinsight}[2][]{
  enhanced, breakable,
  colback=DELightAccent,
  colframe=DEAccent,
  colbacktitle=DEAccent,
  coltitle=white,
  fonttitle=\bfseries\small\sffamily,
  title={Key Insight\quad\textbar\quad #2},
  left=6pt, right=6pt, top=4pt, bottom=6pt,
  arc=4pt, boxrule=0.8pt,
  attach boxed title to top left={yshift=-2mm, xshift=8pt},
  boxed title style={arc=2pt, colback=DEAccent, size=small},
  #1
}

% System Design Note (green)
\newtcolorbox{systemdesign}[2][]{
  enhanced, breakable,
  colback=DELightGreen,
  colframe=DEGreen,
  colbacktitle=DEGreen,
  coltitle=white,
  fonttitle=\bfseries\small\sffamily,
  title={System Design Note\quad\textbar\quad #2},
  left=6pt, right=6pt, top=4pt, bottom=6pt,
  arc=4pt, boxrule=0.8pt,
  attach boxed title to top left={yshift=-2mm, xshift=8pt},
  boxed title style={arc=2pt, colback=DEGreen, size=small},
  #1
}

% Common Misconception (red)
\newtcolorbox{caution}[2][]{
  enhanced, breakable,
  colback=DELightRed,
  colframe=DERed,
  colbacktitle=DERed,
  coltitle=white,
  fonttitle=\bfseries\small\sffamily,
  title={Common Misconception\quad\textbar\quad #2},
  left=6pt, right=6pt, top=4pt, bottom=6pt,
  arc=4pt, boxrule=0.8pt,
  attach boxed title to top left={yshift=-2mm, xshift=8pt},
  boxed title style={arc=2pt, colback=DERed, size=small},
  #1
}

% Production Lesson (gray)
\newtcolorbox{techdebt}[2][]{
  enhanced, breakable,
  colback=DELightGray,
  colframe=DEGray,
  colbacktitle=DEGray,
  coltitle=white,
  fonttitle=\bfseries\small\sffamily,
  title={Production Lesson\quad\textbar\quad #2},
  left=6pt, right=6pt, top=4pt, bottom=6pt,
  arc=4pt, boxrule=0.8pt,
  attach boxed title to top left={yshift=-2mm, xshift=8pt},
  boxed title style={arc=2pt, colback=DEGray, size=small},
  #1
}

% Formal Definition (white + blue border)
\newtcolorbox{defbox}[2][]{
  enhanced, breakable,
  colback=white,
  colframe=DEBlue!55,
  colbacktitle=DEBlue!10,
  coltitle=DEBlue,
  fonttitle=\bfseries\small\sffamily,
  title={Definition\quad\textbar\quad #2},
  left=6pt, right=6pt, top=4pt, bottom=6pt,
  arc=4pt, boxrule=0.8pt,
  #1
}

% Case Study (purple)
\newtcolorbox{casestudy}[2][]{
  enhanced, breakable,
  colback=DELightPurple,
  colframe=DEPurple,
  colbacktitle=DEPurple,
  coltitle=white,
  fonttitle=\bfseries\small\sffamily,
  title={Case Study\quad\textbar\quad #2},
  left=6pt, right=6pt, top=4pt, bottom=6pt,
  arc=4pt, boxrule=0.8pt,
  attach boxed title to top left={yshift=-2mm, xshift=8pt},
  boxed title style={arc=2pt, colback=DEPurple, size=small},
  #1
}

% Chapter Learning Objectives
\newtcolorbox{chapterlearning}[1][]{
  enhanced, breakable,
  colback=DEBlue!5,
  colframe=DEBlue,
  colbacktitle=DEBlue,
  coltitle=white,
  fonttitle=\bfseries\normalsize\sffamily,
  title={Learning Objectives},
  left=8pt, right=8pt, top=6pt, bottom=8pt,
  arc=4pt, boxrule=1pt,
  #1
}

%----------------------------------------------------------------------
% 9. CHAPTER AND SECTION HEADING STYLES
%----------------------------------------------------------------------
\usepackage{titlesec}

% Chapter: oversized gray number, colored title, horizontal rule
\titleformat{\chapter}[display]
  {\normalfont\sffamily}
  {\raggedleft\color{DEBlue!25}\fontsize{90}{90}\selectfont\thechapter}
  {-24pt}
  {\raggedright\color{DEBlue}\fontsize{22}{26}\selectfont\bfseries}
  [\vspace{4pt}\color{DEBlue}\rule{\linewidth}{1.8pt}\vspace{4pt}]

\titlespacing*{\chapter}{0pt}{-40pt}{28pt}

\titleformat{\section}
  {\normalfont\sffamily\Large\bfseries\color{DEBlue}}
  {\thesection}{0.9em}{}
  [\vspace{-4pt}\color{DEBlue!25}\rule{\linewidth}{0.5pt}]

\titleformat{\subsection}
  {\normalfont\sffamily\large\bfseries\color{DEBlue!85}}
  {\thesubsection}{0.9em}{}

\titleformat{\subsubsection}
  {\normalfont\sffamily\normalsize\bfseries\color{DEBlue!75}}
  {\thesubsubsection}{0.9em}{}

%----------------------------------------------------------------------
% 10. HEADERS AND FOOTERS
%----------------------------------------------------------------------
\usepackage{fancyhdr}
\pagestyle{fancy}
\fancyhf{}
\renewcommand{\headrulewidth}{0.4pt}
\renewcommand{\footrulewidth}{0pt}
\fancyhead[LE]{\small\sffamily\color{DEGray}\nouppercase{\leftmark}}
\fancyhead[RO]{\small\sffamily\color{DEGray}\nouppercase{\rightmark}}
\fancyfoot[LE,RO]{\small\sffamily\color{DEGray}\thepage}

\fancypagestyle{plain}{
  \fancyhf{}
  \renewcommand{\headrulewidth}{0pt}
  \fancyfoot[C]{\small\sffamily\color{DEGray}\thepage}
}

%----------------------------------------------------------------------
% 11. CAPTIONS
%----------------------------------------------------------------------
\usepackage{caption}
\usepackage{subcaption}
\captionsetup{
  font=small,
  labelfont={bf,color=DEBlue},
  margin=10pt,
  justification=justified,
}

%----------------------------------------------------------------------
% 12. LISTS
%----------------------------------------------------------------------
\usepackage{enumitem}
\setlist[itemize]{leftmargin=*, topsep=4pt, itemsep=2pt, parsep=0pt}
\setlist[enumerate]{leftmargin=*, topsep=4pt, itemsep=2pt, parsep=0pt}
\setlist[description]{leftmargin=0pt, topsep=4pt, itemsep=2pt}

%----------------------------------------------------------------------
% 13. EPIGRAPH
%----------------------------------------------------------------------
\usepackage{epigraph}
\setlength{\epigraphwidth}{0.62\textwidth}
\renewcommand{\epigraphflush}{flushright}
\renewcommand{\epigraphrule}{0.4pt}

%----------------------------------------------------------------------
% 14. BIBLIOGRAPHY
%----------------------------------------------------------------------
\usepackage[
  style=authoryear,
  backend=biber,
  maxcitenames=3,
  maxbibnames=99,
  giveninits=true,
  uniquename=init,
  dashed=false,
  isbn=false,
  doi=true,
  url=false,
  sorting=nyt
]{biblatex}
\addbibresource{bibliography.bib}

%----------------------------------------------------------------------
% 15. CROSS-REFERENCES AND HYPERLINKS
%----------------------------------------------------------------------
\usepackage[
  colorlinks=true,
  linkcolor=DEBlue,
  citecolor=DEGreen,
  urlcolor=DEAccent,
  pdftitle={Data Engineering: Distributed Systems, Scalable Processing, and Modern Architectures},
  pdfauthor={Rokan Uddin Faruqui},
  pdfsubject={Data Engineering Textbook},
  pdfkeywords={big data, hadoop, spark, data engineering, distributed systems},
  bookmarksnumbered=true,
  bookmarksopen=true,
]{hyperref}
\usepackage[capitalise,noabbrev,nameinlink]{cleveref}

%----------------------------------------------------------------------
% 16. INDEX
%----------------------------------------------------------------------
\usepackage{makeidx}
\makeindex

%----------------------------------------------------------------------
% 17. BLANK PAGES BETWEEN CHAPTERS
%----------------------------------------------------------------------
\usepackage{emptypage}

%----------------------------------------------------------------------
% 18. UTILITY COMMANDS
%----------------------------------------------------------------------
% Defined term (bold + italic) — auto-indexed
\newcommand{\term}[1]{\textbf{\textit{#1}}\index{#1}}
% Software tool name (sans-serif) — auto-indexed
\newcommand{\tool}[1]{\textsf{#1}\index{#1}}
% Inline code
\newcommand{\code}[1]{\texttt{#1}}
% File/dataset name
\newcommand{\dataset}[1]{\texttt{#1}}
% Data-size units (upright, non-breaking)
\newcommand{\KB}{\,\text{KB}}
\newcommand{\MB}{\,\text{MB}}
\newcommand{\GB}{\,\text{GB}}
\newcommand{\TB}{\,\text{TB}}
\newcommand{\PB}{\,\text{PB}}
\newcommand{\EB}{\,\text{EB}}
\newcommand{\ZB}{\,\text{ZB}}

%----------------------------------------------------------------------
% 19. EXERCISE COUNTERS AND MACROS
%----------------------------------------------------------------------
\newcounter{sqctr}[chapter]
\newcounter{apctr}[chapter]
\newcounter{dpctr}[chapter]
\newcounter{pectr}[chapter]

\newcommand{\sq}{\stepcounter{sqctr}%
  \par\vspace{4pt}\noindent
  \textbf{SQ\thechapter.\thesqctr.}\enspace}
\newcommand{\ap}{\stepcounter{apctr}%
  \par\vspace{6pt}\noindent
  \textbf{AP\thechapter.\theapctr.}\enspace}
\newcommand{\dpr}{\stepcounter{dpctr}%
  \par\vspace{6pt}\noindent
  \textbf{DP\thechapter.\thedpctr.}\enspace}
\newcommand{\pe}{\stepcounter{pectr}%
  \par\vspace{6pt}\noindent
  \textbf{PE\thechapter.\thepectr.}\enspace}

% Section heading shorthand for exercise sections
\newcommand{\exercisesection}[1]{%
  \vspace{10pt}%
  {\sffamily\large\bfseries\color{DEBlue} #1}%
  \par\vspace{4pt}%
  {\color{DEBlue!25}\rule{\linewidth}{0.4pt}}%
  \vspace{4pt}%
}
